% c05 —— §4 Complex Fermions（复费米子）与 §5 The Majorana Chain（马约拉纳链）
% 源：QuantumWorkTry.tex 第 1987--2479 行（arXiv:1508.04715v2）
% 说明：全文逐段中译；公式、\label/\ref/\cite 与原文逐字一致；Tutor Note 批注不参与编译。

\section{复费米子（Complex Fermions）}\label{complex}

\begin{keytakeaways}
\begin{itemize}[leftmargin=1.6em,itemsep=0.2em,topsep=0.2em]
  \item \textbf{物理图像}：在不可定向时空上，3+1 维拓扑超导体的边界费米子按 Pin 群的复表示变换，配分函数取复值；单个边界的路径积分带有整体反常，只有与体自旋拓扑量子场论的相位相乘后才在物理上自洽。
  \item \textbf{数学工具}：映射环面上的 APS 不变量、Dai--Freed 定理（普法夫量线丛的截面）、pin 配边不变量，以及 16 次单位根。
  \item \textbf{核心结论}：配分函数相位一般是 16 次单位根，故相互作用下能带数只是模 16 意义下的拓扑不变量；能带数为 8 的倍数时体配分函数化为欧拉示性数的符号，映射环面上欧拉示性数为零使其传统整体反常消失，但要处处定义配分函数仍需能带数为 16 的倍数。
\end{itemize}
\end{keytakeaways}

\subsection{概述（Overview）}\label{overv}

我们关于复费米子的基本例子，是 3 维空间或 4 维时空中的拓扑绝缘体或拓扑超导体。我们在第 \ref{pseudo} 节已经处理过这个问题，但那里只考虑了可定向的时空。这里我们讨论当时空被允许不可定向时会发生什么。在这种情况下\footnote{我们考虑满足 $\sT^2=(-1)^F$ 的系统，因而相应的群是 $\Pin^+(3)$ 而不是 $\Pin^-(3)$。见附录 \ref{unor}。}，边界费米子按 $\Pin^+(3)$ 的复表示变换，而不是按 $\Spin(3) \cong \SU(2)$ 的赝实表示变换。即便只在形式层面上，它们的配分函数现在也是复值的、而非实值的，于是我们必须考虑复费米子的复值路径积分中的反常。

% [Tutor Note: 关键区别在于表示的“实性”：可定向 3-流形上边界费米子取 $\Spin(3)\cong\SU(2)$ 的赝实（pseudoreal）二维表示，配分函数原则上可实；不可定向时取 $\Pin^+(3)$ 的复表示，$\psi$ 与 $\psi'$ 互为复共轭，单个配分函数必然复值。$\Pin^+(3)$ 的两个二维不可约复表示互为复共轭，详见附录 \ref{wwod}。原文第 1995 行 pseudreal 为 pseudoreal 的拼写笔误。]

我们首先讨论拓扑超导体，事实证明它是更有趣的情形。讨论以多少熟悉的方式展开。我们先考虑边界模的反常，这些边界模是一个 $d=3$ 的马约拉纳费米子。一如既往，该反常可以表示为映射环面的一个拓扑不变量。对于一般的复费米子 \cite{Global}，映射环面的相关不变量是 $\exp(-\i\pi\eta/2)$，其中 $\eta$ 是高一维的某个狄拉克算子的 APS $\eta$ 不变量——在我们的例子里，就是一个 $D=4$ 马约拉纳费米子的狄拉克算子。接着我们考虑与体拓扑超导体相伴的 sTQFT 的配分函数。利用 $D=4$ 中拓扑超导相变的一个标准刻画，我们证明体 sTQFT 的配分函数同样是 $\exp(-\i\pi\eta/2)$。与我们此前作为 sTQFT 配分函数遇到的其它体不变量一样，这个量无法用局域的、保持对称性的边界条件来定义。我们可以借助 APS 边界条件以保持对称性的方式来定义它，但那样它就不是拓扑不变量，单凭它自身仍然不能令人满意。不过，乘积 $|Z_\psi|\exp(-\i\pi \eta/2)$ 在物理上是合理的，Dai--Freed 定理 \cite{DF} 说明了这一点。该乘积应被理解为边界费米子与体 sTQFT 合并后的配分函数。

% [Tutor Note: 此处的归一化极易混淆：复费米子的映射环面不变量是 $\exp(-\i\pi\eta/2)$，其中 $\eta$ 是高一维（此处 $D=4$）狄拉克算子的 APS $\eta$ 不变量；而对实/赝实费米子，前文的不变量写作 $(-1)^{\I/2}$。两者在可定向情形一致（见第 \ref{bpf} 节的核对），因为可定向时 $(-1)^{\I/2}=\exp(-\i\pi\eta/2)$。指数上出现 $\eta/2$ 而非 $\eta$，反映的是普法夫量相对于行列式只取“一半自由度”。]

最后，我们考虑一个在此前本文所研究的情形中没有近似对应的问题。它涉及一个 $D=4$ 系统，其边界上承载 $\nu$ 个相同的无能隙马约拉纳费米子。在自由费米子层面，$\nu$ 是取整数值的不变量；但已有论证 \cite{FCV,WS,MFCV,KitTwo} 表明，在相互作用层面 $\nu$ 只是 mod 16 意义下的不变量。从目前的观点看，这一点成立是因为：在 $D=4$ 中 $\exp(-\i\pi\eta/2)$ 一般是 1 的 16 次根，而在带空间边界的 4-流形上，除非 $\nu$ 是 16 的倍数，否则 $\exp(-\nu\cdot \i\pi\eta/2)$ 没有令人满意的、保持对称性的定义。

% [Tutor Note: “为何是 16”的路径积分理由：单个 sTQFT 的配分函数是 $\exp(-\i\pi\eta/2)$；$\nu$ 个副本要求把该相位取 $\nu$ 次幂，而该相位本身只有 16 个可能值，故只有 $\nu$ 模 16 可观测。这个 16 来自 $\eta$ 在 $D=4$ 中的 $1/4$ 量子化（附录 \ref{feta}）；对 $\Pin^-$ 的 $D=2$ 情形则是 $1/2$ 量子化，给出 8（见第 \ref{maj} 节的马约拉纳链）。]

关于 $\nu=16$ 这一数值，\cite{KTTW} 曾基于如下设想给出一个多少类似的解释：相关的 sTQFT 配分函数应当是一个配边不变量。一个显式的整体反常计算 \cite{HCR} 表明，含 $\nu$ 个相同马约拉纳费米子的边界理论除非 $\nu$ 是 8 的整数倍，否则就有反常。$\nu=8$ 时的情形其实相当微妙。我们将证明：尽管 $\exp(-\i\pi\eta/2)$ 一般是 1 的 16 次根，在映射环面的情形它却是 1 的 8 次根。因此 $\nu=8$ 的理论没有可通过映射环面构造探测到的传统反常。然而，$\nu=8$ 时还存在一个更微妙的问题，它类似于我们在第 \ref{signp} 节针对仅在可定向 3-流形上考虑单个马约拉纳费米子所解释的情形：尽管不存在传统意义上的反常，也同样没有自然的方式去指定配分函数的符号。在第 \ref{signp} 节，原因在于 $(-1)^{\I/2}$ 一般是非平凡的，尽管对映射环面它等于 1；在目前语境下，原因则是类似的说法对不变量 $\exp(-8\cdot\i\pi\eta/2)$ 成立。

% [Tutor Note: 这里给出“8 而非 16 在映射环面层面出现”的机制：一般 $\Pin^+$ 4-流形上 $\exp(-\i\pi\eta/2)$ 是 16 次根，故只有 $\nu\in 16\Z$ 才能让 $\nu$ 个副本的 sTQFT 平凡；但映射环面 $X\to S^1$ 上 $\eta$ 的取值被限制，$\exp(-\i\pi\eta/2)$ 落入 8 次根的子群，于是 $\nu\in 8\Z$ 就足以消除该构造能探测到的相位（第 \ref{zeb} 节用 $\chi(X)=0$ 说明这一点）。更精细的判据要看所有 4-流形而非仅映射环面。]

$\nu=8$ 的问题还有一个有趣性质，足以显示其微妙之处。对 $\nu=8$ 理论——也就是 $\nu=1$ 的 sTQFT 的 8 次幂——实际上可以定义一个无拓扑序的、有能隙的、保持对称性的边界态；但这个边界态与幺正性不相容。

% [Tutor Note: 这是 $\nu=8$ 特殊性的最深刻之处：单看切割--黏合公理，体理论 $(-1)^{w_4}$ 允许一个对称保持的有隙边界态；但反射正定性（幺正性）会排除它。第 \ref{zeb} 节末尾用四维球体 $\chi=1$ 的黏合论证给出具体反例，请对照阅读。]

$d=3$ 拓扑绝缘体也通过同样的论证与 $\eta$ 不变量相联系。但在这种情形下，不存在整数值参数 $\nu$ 的类似物。因此结论不那么有趣，而且走向不可定向时空并不会揭示出超出第 \ref{pseudo} 节所探讨内容之外的结构。

\subsection{边界费米子的反常（Anomaly Of The Boundary Fermions）}\label{anombo}

不可定向的 3-流形 $Y$ 总可以构造为商 $\h Y/\Z_2$，其中 $\h Y$ 是 $Y$ 的可定向双重覆盖，而 $\Z_2$ 由满足 $\tau^2=1$ 的定向反转对称 $\tau:\h Y\to \h Y$ 生成。设给定了 $\h Y$ 上一个 $\tau$-不变的自旋结构 $\S$。这特别意味着：我们有 $\h Y$ 上费米子场 $\h\psi$ 的概念，以及 $\tau$ 对 $\h\psi$ 的作用，且满足 $\tau^2=1$。

在 $Y$ 上，我们可以通过如下方式定义一个 $\pin^+$ 结构 $\P$：带有此 $\pin^+$ 结构的 $Y$ 上费米子场 $\psi$，就是 $\h Y$ 上满足 $\tau\h\psi=\h\psi$ 的费米子场 $\h\psi$。此外，$Y$ 上每一个 $\pin^+$ 结构都以这种方式出现，对应于 $Y$ 上某个 $\tau$-不变的自旋结构 $\S$ 的选取。

% [Tutor Note: 这段建立了不可定向 $Y$ 上 $\pin^+$ 结构与 $\h Y$ 上 $\tau$-不变自旋结构之间的对应：$Y$ 上的 $\pin^+$ 旋量即 $\h Y$ 上满足 $\tau\h\psi=\h\psi$ 的旋量。原文第 2056 行 “for some choice of $\tau$-invariant spin structure $\S$ on $Y$” 中 “on $Y$” 疑为 “on $\h Y$” 之误，因为 $\S$ 按定义是 $\h Y$ 上的自旋结构；译文保留原文以便对照。]

但我们也可以在 $Y$ 上定义一个“互补的”$\pin^+$ 结构 $\P'$：带有这个替代 $\pin^+$ 结构的 $Y$ 上费米子场 $\psi'$，就是 $\h Y$ 上满足 $\tau\h\psi=-\h\psi$ 的费米子场 $\h\psi$。事实上，$\P$ 与 $\P'$ 之间的关系是完全对称的，两者没有哪一个更优越。当给定几何对称 $\tau:\h Y\to \h Y$ 以及一个在自同构意义下 $\tau$-不变的自旋结构 $\S\to \h Y$ 时，这并不能告诉我们应当让 $\tau$ 以何种符号作用在 $\S$ 上。符号的选择决定了我们所谓的 $\P$ 与所谓的 $\P'$ 各是什么。

在 $\h Y$ 上，我们可以定义通常的厄米狄拉克算子 $\D=\i\slashed{D}=\i\sum_k\g^k D_k$。然而在奇数维中，狄拉克算子与定向反转对称反对易。这一点已在式（\ref{refact}）中说明过，并且正如那里所解释的，这正是 $2+1$ 维狄拉克方程中若加入质量项便会破坏反射对称的原因。

% [Tutor Note: 奇维中 $\D$ 在定向反转下为奇，是因为定向反转除变换坐标外还要作用在伽马矩阵上（$\g^k\to-\g^k$），从而使 $\slashed D$ 变号；式（\ref{refact}）给出同一事实的坐标表述。由此 $2+1$ 维质量项 $m\bar\psi\psi$ 在反射下变号，故在有反射对称的边界上被禁戒——这正是“拓扑绝缘体边界无能隙”的场论根源。]

特别地，$\D$ 在 $\tau$ 下为奇，因此若 $\tau\psi=\pm\psi$，则 $\tau(\D\psi)=\mp\D\psi$。在 $Y$ 上解读，这句话意味着：不存在作用在作为 $\P$ 的截面的费米子场 $\psi$ 上、或作用在作为 $\P'$ 的截面的费米子场 $\psi'$ 上的自伴狄拉克算子。相反，$Y$ 上自然的狄拉克算子把 $\P$ 映到 $\P'$、把 $\P'$ 映到 $\P$。算子 $\D_\P:\P\to \P'$ 与 $\D_{\P'}:\P'\to \P$ 互为伴随，于是我们可以把它们合成一个自伴算子；在基 $\bp \P\cr \P'\ep$ 下它是
\be\label{intr}\D=\bp 0 & \D_{\P'}\cr \D_\P & 0 \ep. \ee
事实上，$\D$ 就是 $\h Y$ 上通常的自伴狄拉克算子（对于自旋结构 $\S$），作用在被重新解释为 $Y$ 上的一对场 $\psi, $ $\psi'$ 的旋量场 $\h\psi$ 上。

% [Tutor Note: 容易忽略的要点：$\P$ 与 $\P'$ 各自都不是自伴算子能作用的空间，$\D_\P$ 与 $\D_{\P'}$ 只能成对出现；式（\ref{intr}）的块结构正是把 $\h Y$ 上的自伴算子用 $Y$ 上的数据写出来。$\D_\P$ 与 $\D_{\P'}$ 互为伴随这一点依赖于 $\h Y$ 上狄拉克算子的自伴性与 $\tau$ 的等距性，矩阵中上下两个非对角块的位置也由此确定，不能随意调换。]

更根本地说，对 $\psi$ 或对 $\psi'$ 不存在自伴的欧氏号差狄拉克算子的原因在于：在欧氏号差下它们是复费米子，而非实或赝实费米子。正如附录 \ref{wwod} 所解释的，群 $\Pin^+(3)$ 有两个不同的二维表示，它们互为复共轭。场 $\psi$ 与 $\psi'$ 分别联系于这两个不同的表示。

尽管我们无法为单独的 $\psi$ 或 $\psi'$ 写出自伴狄拉克算子，但这对仅含 $\psi$（或同样地仅含 $\psi'$）的经典狄拉克作用量
\be\label{terb}\int\d^3x \sqrt g\bar\psi\i\slashed{D}\psi \ee
并无妨碍。回忆一下，正如第 \ref{ferp} 节所讨论的，狄拉克作用量实际涉及的是一个反对称双线性型 $\DD$，而不是一个算子 $\D$。在 $\h Y$ 上，两者之间的关系是 $\DD_{\gamma\beta}=\veps_{\gamma\alpha}\D^\alpha_\beta$，其中 $\veps_{\gamma\beta}$ 是 $\Spin(3)=\SU(2)$ 的二维表示上的反对称双线性型。但该型在 $\tau$ 下为奇（更一般地说，它在 $\Pin^+(3)$ 的非连通分支下为奇）。由于 $\D^\alpha_\beta$ 在 $\tau$ 下也为奇，由此可知 $\DD_{\gamma\beta}$ 为偶，因而只能对 $\tau$-偶的费米子 $\psi$ 或 $\tau$-奇的费米子 $\psi'$ 定义。

% [Tutor Note: 这是全文最易出符号错误的地方之一：$\veps_{\gamma\beta}$（$\SU(2)$ 的不变反对称张量）在 $\tau$ 下为奇，而 $\D^\alpha_\beta$ 在 $\tau$ 下也为奇，两个“奇”相乘才使 $\DD$ 为偶，从而经典作用量才能与 $\tau$ 的投影相容。若误把 $\veps$ 当作偶，就会得出 $\psi$ 与 $\psi'$ 都无法定义作用量的错误结论。原文第 2087 行 “Up on $\h Y$” 系 “Upon $\h Y$” 的笔误。]

我们所说的一切，在比如 2 维的可定向流形上有一个大概熟悉得多的类比。在 2 维中，我们可以考虑手征费米子 $\psi_+$ 与 $\psi_-$，它们在欧氏号差下互为复共轭。2 维狄拉克算子把 $\psi_+$ 映到 $\psi_-$、把 $\psi_-$ 映到 $\psi_+$，所以要写出自伴狄拉克算子就必须把两个手征性合起来。但狄拉克作用量把 $\psi_+$ 与 $\psi_+$ 配对、把 $\psi_-$ 与 $\psi_-$ 配对，因此（在反常抵消的问题之外）考虑只含单一手征性的理论是有意义的——例如 $2+1$ 维手征拓扑超导体的边缘模。

回到可能不可定向的 3-流形上的场 $\psi$ 与 $\psi'$：$\psi'$ 的狄拉克作用量与路径积分是 $\psi$ 的相应量的复共轭，所以两个配分函数 $Z_\psi$ 与 $Z_{\psi'}$ 互为复共轭。事实上，具体地说，$\psi$ 与 $\psi'$ 合起来就是 $\h Y$ 上的费米子 $\h\psi$，因此形式上 $Z_\psi Z_{\psi'}=Z_{\h\psi}$，而我们断言 $Z_{\h\psi}\geq 0$。一如既往，我们有 $Z_{\h\psi}=|Z_{\h\psi}|\exp(-\i\pi\eta/4)$。然而，由于 $\h Y$ 的狄拉克算子 $\D$ 在 $\tau$ 下为奇，其谱在 $\lambda \leftrightarrow -\lambda$ 下对称，所以 $\h Y$ 上的非零模对 $\eta$ 没有贡献。于是要么没有零模，此时 $\eta=0$ 且 $Z_{\h\psi}>0$；要么存在零模，此时 $Z_{\h\psi}=0$。

% [Tutor Note: $\eta$ 只由零模贡献，是因为 $\tau$-奇性保证非零本征值成对 $\lambda,-\lambda$ 出现、贡献相互抵消；而按式（\ref{uv}）的约定 $\sign(0)=+1$，每个零模对 $\eta$ 贡献 $+1$。故无零模时 $\eta=0$，$Z_{\h\psi}$ 严格为正；有零模时普法夫量直接为零。注意这里的 $\eta$ 是 3 维（$\h Y$ 上）算子的 $\eta$，与 4 维 $\eta$ 差一维，指数分母相应为 $4$ 而非 $2$。]

虽然乘积 $Z_\psi Z_{\psi'}$ 为正、从而完全无反常，但 $Z_\psi$ 与 $Z_{\psi'}$ 各自是复值的并且带有反常。它们没有微扰反常（一般地，无边界奇数维时空中没有微扰反常），但有整体反常。确实，$Z_\psi$ 的整体反常最近已在 \cite{HCR} 中就一个具体例子计算过。这与下述更熟悉的事实类似：在 2 维时空中的手征费米子 $\psi_+$ 与 $\psi_-$ 的情形，$Z_{\psi_+}$ 与 $Z_{\psi_-}$ 带有反常（在该情形中既有微扰反常也有整体反常），但乘积 $Z_{\psi_+}Z_{\psi_-}$ 非负且无反常。

一个计算复费米子路径积分中整体反常的一般框架，已在 \cite{Global} 第 4 节中描述。给定 $Y$ 的一个微分同胚 $\phi$（在更一般的问题中，是微分同胚与规范变换的复合），如第 \ref{trada} 节那样构造映射环面 $X$。$X$ 是把 $Y\times I$ 的两端通过 $\phi$ 黏合起来所得的 4-流形，其中 $I$ 是区间 $0\leq s\leq 1$。我们给 $X$ 配以度规
\be\label{dif}\d \ell^2= \d s^2+\epsilon^2 g_s,\ee
其中 $g_s$ 是 $Y$ 上依赖于 $s$ 的一族度规，$\epsilon$ 是适当小的参数。在基（\ref{intr}）下，引入一个新的伽马矩阵
\be\label{mezzo}\g_s=\bp 1 & 0 \cr 0 & -1 \ep. \ee
然后定义 $X$ 上的狄拉克算子：
\be\label{ezzo}\D_X=\i\g_s\frac{D}{D s}+\frac{1}{\eps} \bp 0 & \D_{\P'}\cr \D_\P & 0 \ep. \ee
这里 $\D_Y=\bp 0 & \D_{\P'}\cr \D_\P & 0 \ep$ 是 $Y$ 上带 $s$ 依赖度规 $g_s$ 的狄拉克算子，而显式的因子 $1/\epsilon$ 来自式（\ref{dif}）中我们用的是度规 $\eps^2 g_s$。\cite{Global} 中分析的结果表明，$Z_\psi$ 在 $\phi$ 下的变化可以用 $D=4$ 狄拉克算子 $\D_X$ 的 $\eta$ 不变量表出：
\be\label{wezzo}Z_\psi\to Z_\psi \exp(-\i\pi\eta/2). \ee

% [Tutor Note: 度规（\ref{dif}）在 $s$ 方向取单位长度、在 $Y$ 方向取 $\eps^2 g_s$，所以 $Y$ 方向狄拉克算子的本征值标度为 $1/\eps$；式（\ref{ezzo}）中 $Y$ 方向块矩阵前乘 $1/\eps$ 正是这个原因。$\g_s=\diag(1,-1)$ 是在 $\P\oplus\P'$ 基下沿 $s$ 方向的伽马矩阵，它使 $\D_X$ 成为 4 维自伴算子。原文 “familiy”（第 2120 行）系 “family” 的拼写笔误。]

我们将不详细重复这一计算，只提及几点。由于我们要计算的是一个拓扑不变量，可以取一般的依赖于 $s$ 的度规族 $g_s$，使得 $\D_Y$ 作为 $s$ 的函数没有能级交叉。再取 $\epsilon$ 足够小，就可以约化为如下情形：$\D_X$ 的谱、从而其 $\eta$ 不变量，可以用绝热近似计算。\footnote{若 $\lambda_i(s)$ 是 $\epsilon^{-1}\D_Y$ 作为 $s$ 的函数的本征值（我们假定除来自赝实性的二重简并外它们非简并），则绝热近似的条件是：与相关的本征值差相比，$\d\lambda_i/\d s$ 应足够小。在没有能级交叉的假设下，只要 $\epsilon$ 足够小，该条件即满足。}在绝热近似下，$\D_Y$ 的每个本征值对 $\exp(-\i\pi\eta/2)$ 的贡献本质上来自其贝里联络的和乐。（关于给定本征态绝热演化的贝里联络，在 \cite{Global} 第 215 页中部有所引入，尽管并未使用这个名称。）在绝热近似下，同一个和乐给出每个本征值对 $Z_\psi$ 在 $\phi$ 下相位变化的贡献。

% [Tutor Note: 绝热图像：沿 $s\in S^1$ 一圈，每个本征态获得贝里相位 $\exp(-\i\oint A_i)$，$\exp(-\i\pi\eta/2)$ 的相位即这些贡献之积；同一和乐也解释 $Z_\psi$ 的相位变化，这正是式（\ref{wezzo}）的微观来源。脚注中的本征值指 $\epsilon^{-1}\D_Y$（物理标度后的 $Y$ 方向算子），其能隙随 $1/\eps$ 增大，故 $\eps\to0$ 使绝热条件成立——这是原文省略的“为何小 $\eps$ 可实现绝热近似”的关键一步。]

式（\ref{wezzo}）的一个推论是：正如我们此前研究过的其他例子，反常是一个配边不变量，意思是若 $X$ 是某个 5-流形 $Z$ 的边界、且 $\pin^+$ 结构可延拓到其上，则该反常为平凡。这里我们用到式（\ref{molg}）。在 5 维 $\pin^+$ 流形上，由于赝实性，指标 $\I$ 总是偶数，因此式（\ref{molg}）表明：若 $X$ 是一个边界，则 $\eta$ 是 4 的倍数，从而 $\exp(-\i\pi\eta/2)=1$。

% [Tutor Note: 论证链条：5 维 $\pin^+$ 流形上 $\I$ 为偶数（赝实性/Kramers 型加倍）；式（\ref{molg}）给出 $\I=-\eta/2$，故 $\eta\in 4\Z$。注意 $\eta$ 的绝对值没有拓扑意义，重要的是它模 4 的取值：$\exp(-\i\pi\eta/2)$ 在 $\eta\to\eta+4$ 下不变。]

真正查阅 \cite{Global} 以及后来作出重要改进的 \cite{DF} 的读者会发现，文中假定了 $Y$ 为偶数维、$X$ 为奇数维。原因很简单：若 $Y$ 可定向，则 $Y$ 上的复费米子只在其维数为偶数时才存在。当时并未设想涉及不可定向流形的应用，因此假定 $Y$ 为偶数维是自然的。然而，这两篇文章中的推理对任意维数（偶数或奇数）的复费米子同样适用。

就结果的应用而言，一个重要差别来自这样一个事实：在偶数维中，复费米子一般会产生微扰反常。\footnote{在 \cite{Global} 中（而非 \cite{DF} 中），某些地方假定了微扰反常相消。}此外，若 $Y$ 是偶数维，则 $X$ 是奇数维；而若它是某个 $Z$ 的边界，则 $Z$ 又是偶数维。若 $\D_X$ 是 $X$ 上的狄拉克算子、$\eta$ 是其 $\eta$ 不变量，则典型地 $\exp(-\i\pi\eta/2)$ 不是配边不变量（甚至不是拓扑不变量），因为 $Z$ 上的 APS 指标定理中存在曲率项。（对 4 维的 $Z$，该曲率项就是式（\ref{trof}）中的 $\hA-P$ 项。）然而，在偶数维流形 $Y$ 上，对于微扰反常相消的复费米子，配边不变性会重新成立。\footnote{一个例子是 $D=4$ 中规范群为 $G=\SU(3)\times \SU(2)\times \U(1)$ 的粒子物理标准模型。在欧氏号差下，费米子处于 $\Spin(4)\times G$ 的一个复但高度可约的表示中。该表示使得微扰反常相消。}这可由 APS 指标定理给出的 $\eta$ 公式推出。当然，这里的整体反常来自一个 $\eta$ 不变量，而它在给定理论中接收所有费米子多重态的贡献，因此在运用 APS 指标定理时，我们要对所有多重态求和。$Z$ 上 APS 指标定理中的曲率项正是费米子的反常多项式，所以在对多重态求和之后，它在无微扰反常的理论中相消。于是 APS 指标定理表明：在任何没有微扰反常的理论中，无论 $D$ 为偶数还是奇数，控制整体反常的不变量 $\exp(-\i\pi\eta/2)$ 都是配边不变量。

% [Tutor Note: 关键机制：APS 指标定理把 $Z$ 上的 $\eta$ 分解为陈--西蒙斯项加曲率（$\hA$）项；曲率项恰是费米子的反常多项式。只有当所有多重态的反常多项式之和为零（无微扰反常）时，$\exp(-\i\pi\eta/2)$ 才只依赖于 $X$ 的配边类。标准模型是这一机制的范例：其表示虽复杂，却恰好使微扰反常相消，故其整体反常是配边不变量。]

$\eta$ 与整体反常的关系，是一个远为著名的结果的近亲或轻微改进：$D$ 维的微扰反常与 $D+1$ 维的陈--西蒙斯项以及 $D+2$ 维的曲率多项式相关。事实上，APS 指标定理蕴含着（见式（\ref{moj}）--（\ref{oj}））$\eta$ 与陈--西蒙斯项之间的密切关系。通常，两者之间的区别只在人们提出相当微妙的问题（例如关于整体反常的问题）时才重要。

\subsection{体配分函数（Bulk Partition Function）}\label{bpf}

这里我们将计算 $3+1$ 维拓扑超导体的体配分函数。其策略应当从第 \ref{why} 节与第 \ref{bulkboundary} 节起就已熟悉。

$3+1$ 维中普通超导体与拓扑超导体之间相变的一个标准描述是：当 $D=4$ 马约拉纳费米子 $\psi$ 的质量 $m$ 穿过 0 时发生相变。即便在不可定向的 4-流形 $X$ 上，也可以定义一个厄米狄拉克算子 $\D$（我们记得第 \ref{anombo} 节指出，在 3 维时空中这并不成立），并且由于马约拉纳费米子按 $\Pin^+(4)$ 的赝实表示变换，其谱被 Kramers 型加倍机制加倍。关于这些陈述的细节，见附录 \ref{do} 与 \ref{fwod}。马约拉纳费米子的路径积分是一个普法夫量，可以通过从每一对中取一个本征值来计算。我们用 $m$ 表示 $\psi$ 的裸质量，并引入一个质量为 $\mu$、满足同一狄拉克方程但具有相反统计的泡利--维拉尔斯场来正规化配分函数 $Z_\psi$。正规化后的配分函数为
\be\label{regpar}Z_{\psi,\reg}=\prod_k{}'\,\frac{\lambda_k+\i m}{\lambda_k+\i\mu},  \ee
其中一如既往，$\prod_k'$ 是“从每一对中只取一个本征值”的指令。

% [Tutor Note: 为什么是普法夫量而非行列式：赝实表示使每个本征值至少二重简并（Kramers 型），$\prod_k'$ 只取每对中的一个，故得到的是 $\Pf(\D)$ 而非 $\det(\D)$；这也解释了后面 $\upeta\to\eta/2$ 的因子。泡利--维拉尔斯场（质量 $\mu$、相反统计）的作用是抵消紫外发散；关键物理只依赖 $m$ 与 $\mu$ 的相对符号，因为在大体积极限下可连续调节而穿过相变点。]

本质上只有两种情形，取决于 $m$ 与 $\mu$ 同号还是反号。若 $m$ 与 $\mu$ 同号，我们得到拓扑平庸的超导体。确实，在 $X$ 非常大的极限下，$m$ 与 $\mu$ 的具体数值无关紧要，只有它们的符号有意义。若 $m$ 与 $\mu$ 同号，我们不妨取 $m=\mu$，此时 $Z_{\psi,\reg}=1$，完全平庸。

有趣的情形是 $m$ 与 $\mu$ 反号。此时不妨取 $m=-\mu$，于是
\be\label{negpar}Z_{\psi,\reg}=\prod_k{}'\,\frac{\lambda_k-\i\mu}{\lambda_k+\i\mu}. \ee
分析这一表达式的一个捷径是把它因式分解：
\be\label{wegpar}Z_{\psi,\reg}=\prod_k{}'\,\frac{\lambda_k}{\lambda_k+\i\mu}\cdot \prod_k{}'\,\frac{\lambda_k-\i\mu}{\lambda_k} . \ee
第二个因子是第一个因子的复共轭之逆，所以它们贡献相同的相位。第一个因子已在第 \ref{zelfan} 节分析过，结果由式（\ref{rv}）给出。使用这一结果时，我们必须把 $\upeta$（仅在正电荷费米子空间中计算）替换为 $\eta/2$（其中 $\eta$ 用狄拉克算子的全部本征值计算，而因子 1/2 反映了 $\prod_k'$ 符号）。熟悉式（\ref{riv}）与（\ref{roov}）的读者对这一替换 $\upeta\to \eta/2$ 应当并不陌生。于是式（\ref{wegpar}）中的两个因子各贡献一个因子 $\exp(-\i\pi\eta/4)$，而拓扑超导体的体 sTQFT 配分函数的最终结果是 $\exp(-\i\pi\eta/2)$。

% [Tutor Note: 这里解释了 $\nu$ 个能带情形中指数上的 $\nu/2$ 从何而来：每个“能带”（单个马约拉纳费米子连同其正规化子）贡献 $\exp(-\i\pi\eta/2)$，其中 $\eta/2$ 的因子来自 $\prod_k'$（只数一半本征值）；$\nu$ 个副本相乘即得 $\exp(-\nu\i\pi\eta/2)$。$\upeta\to\eta/2$ 的替换与式（\ref{riv}）、（\ref{roov}）中的约定一致（$\upeta$ 定义在正电荷子空间、$\eta$ 遍及全部本征值）。]

作为检验，让我们验证：若 $X$ 可定向，这一答案与第 \ref{signp} 节所断言的结果 $(-1)^{\I/2}$ 一致。若 $X$ 可定向，我们可以定义手征算子 $\bar\g$（式（\ref{chiro}）），而操作 $\psi\to\bg\psi$ 确保 $\D$ 的非零本征值在 $\lambda\leftrightarrow-\lambda$ 下配对。（由于赝实性，谱对每个 $\lambda$ 还具有偶重数，所以非零模组成四重态 $\lambda,\lambda,-\lambda,-\lambda.$）配对 $\lambda,-\lambda$ 对 $\eta$ 没有贡献，因此对可定向的 $X$，$\eta$ 只接收零模的贡献。设 $n_+$ 与 $n_-$ 分别是正、负手征零模的数目。每个零模对 $\eta$ 贡献 $+1$（见式（\ref{uv}）），所以 $\eta=n_++n_-$ 且 $\exp(-\i\pi\eta/2)=\exp(-\i\pi (n_++n_-)/2)$。另一方面 $\I=n_+-n_-$，所以 $(-1)^{\I/2}=(-1)^{(n_+-n_-)/2}$。由于 $n_+$ 与 $n_-$ 为偶数，两者相等。若要更直接的说明，读者可以回到式（\ref{negpar}）看出：四重态本征值 $\lambda,\lambda,-\lambda,-\lambda$ 没有贡献，而一对零模贡献一个因子 $-1$，正与 $(-1)^{\I/2}$ 中的情形相同。

% [Tutor Note: 核对细节：$n_\pm$ 均为偶数来自赝实性（零模成对），因此 $\exp(-\i\pi(n_++n_-)/2)$ 与 $(-1)^{(n_+-n_-)/2}$ 只差因子 $(-1)^{n_-}=1$。注意式（\ref{uv}）约定 $\sign(0)=+1$，故零模对 $\eta$ 的贡献为正；这也是为什么 $\eta=n_++n_-$ 而非 $n_+-n_-$。这个 $+1$ 约定是后面 $\eta$ 的模 4 量子化的关键。]

注意在论证过程中我们已经证明：若 $X$ 可定向，则 $\eta$ 是偶整数。但若 $X$ 可定向且具有定向反转对称，则 $n_+=n_-$，从而 $\eta$ 是 4 的倍数。

\subsection{反常抵消与 Dai--Freed 定理（Anomaly Cancellation And The Dai-Freed Theorem）}\label{comb}

一如既往，在带边界的 4-流形 $X$ 上定义体 sTQFT 的配分函数 $\exp(-\i\pi\eta/2)$ 存在一个问题。对于欧氏号差下 4 维狄拉克算子，不存在可用于给出令人满意的 $\eta$ 定义的、局域的、保持对称性的自伴边界条件。我们可以用非局域的 APS 边界条件以保持对称性的方式定义 $\eta$。但这样定义的 $\exp(-\i\pi\eta/2)$ 不是拓扑不变量，且只有与边界上的无能隙模结合时才有物理意义。

在第 \ref{anombo} 节我们解释过，拓扑超导体的通常边界费米子是复费米子，当在不可定向的 3-流形 $Y$ 上表述时其路径积分是复值的。此外，按照复费米子路径积分的一般整体反常公式 \cite{Global}，该反常涉及 $\exp(-\i\pi\eta/2)$，其中如上一段所述，$\eta$ 是一个 $D=4$ 马约拉纳费米子的 $\eta$ 不变量。

正如我们研究过的其他情形，在物理上合理的是乘积
\be\label{durf}|\Pf(\DD)|\exp(-\i\pi\eta/2), \ee
其中 $\Pf(\DD)$ 是边界费米子的普法夫量，$\eta$ 是用 APS 边界条件计算的体 $\eta$ 不变量。事实上，“这一乘积在物理上合理”这一断言，本质上就是我们所谓的 Dai--Freed 定理 \cite{DF} 的内容，它是一般整体反常公式的一个改进。

% [Tutor Note: Dai--Freed 定理的精确表述：$\exp(\i\pi\eta/2)$ 自然地是普法夫量线丛 $\PF(\D)$ 的一个截面；等价地，$\Pf(\DD)\exp(-\i\pi\eta/2)$（两因子分别是互为逆的线丛的截面）作为复数良定义。注意两个因子的线丛方向相反：若 $\exp(\i\pi\eta/2)$ 属于 $\PF(\D)$，则 $\exp(-\i\pi\eta/2)$ 属于 $\PF(\D)^{-1}$，二者相乘才能抵消。原文此处写作 $|\Pf(\DD)|$ 而非 $\Pf(\DD)$，是因为在边界算子无零模的区域可取自然平凡化使普法夫量为正（见下一段）。]

我们不打算描述 Dai--Freed 定理的证明，但会作几点或许有用的评论。首先，该定理在 \cite{DF} 中的表述其实更为抽象。为解释这一点，我们应当说明：关于反常的通常数学观点是，尽管 $\Pf(\DD)$ 不能被自然地定义为一个复数、或定义为 $Y$ 上度规的复值函数，但总存在 $Y$ 上度规空间（在更一般情形中是度规加规范场）上的一个酉复线丛，使得 $\Pf(\DD)$ 作为该线丛的截面被典范地定义。这个线丛称为普法夫量线丛，我们将把它记为 $\PF(\DD)$。\cite{DF} 中陈述的 Dai--Freed 定理断言：$\exp(\i\pi\eta/2)$ 作为 $\PF(\D)$ 的截面被自然地定义。等价的陈述是：乘积 $\Pf(\DD)\exp(-\i\pi\eta/2)$（其中两个因子是互为逆的线丛的截面）作为复数良定义。然而，只要我们限制在边界狄拉克算子无零模的轨迹上，线丛 $\PF(\DD)$ 就有自然的平凡化。相对于这一平凡化，普法夫量 $\Pf(\DD)$ 为正，可以替换为其绝对值，而 $\exp(-\i\pi\eta/2)$ 同样是一个良定义的复数，\footnote{当边界狄拉克算子确实有零模时，我们无需讨论 $\exp(-\i\pi\eta/2)$ 的含义，因为此时 $\Pf(\DD)=0$。正如 \cite{DF} 所解释的，这种情形很微妙，因为当边界费米子有零模时，APS 边界条件的定义本身就很微妙。}并采用通常的定义（\ref{roov}）。

最后，我们应当具体解释“乘积 $|\Pf(\DD)|\exp(-\i\pi\eta/2)$ 在物理上合理”这句话的含义。对于实或赝实费米子，一个要点是：当边界狄拉克算子出现零本征值时，$|\Pf(\DD)|(-1)^\zeta$ 或 $|\Pf(\DD)|(-1)^{\I/2}$ 的行为是合理的。之所以如此，是因为每当边界费米子的路径积分按理应当变号时，$(-1)^\zeta$ 或 $(-1)^{\I/2}$ 的符号恰好跳变。对于复费米子，我们处理的不是符号变化，因为 $\Pf(\DD)$ 并非自然地取实值。但弄清零模出现时会发生什么仍然至关重要。

在边界狄拉克算子出现一对零模的位置附近，相应的双线性型 $\DD$ 在其标准形下具有一个 $2\times 2$ 块
\be\label{zumo}\bp 0 & \lambda \cr -\lambda & 0\ep, \ee
其中 $\lambda$ 很小。这里 $\lambda$ 是复数，要使它为零，我们必须变动 $Y$ 上度规和/或规范场中的两个参数。例如，可能存在两个参数 $u_1,u_2$ 使得
\be\label{umo}\lambda=u_1+\i u_2. \ee
等价地，令 $u_1+\i u_2=\varrho e^{\i\varphi}$，我们有 $\lambda=\varrho e^{\i\varphi}$。我们预期当 $\lambda\to 0$ 时路径积分正比于 $\lambda$。因此，为使组合的体与边界配分函数公式 $|\Pf(\DD)|\exp(-\i\pi\eta/2)$ 在 $\lambda=0$ 附近物理上合理，$\exp(-\i\pi\eta/2)$ 在 $\lambda=0$ 附近必须不光滑。相反，它必须正比于 $e^{\i\varphi}$，而后者在 $\lambda=0$ 处没有良好定义。具体地说，这是因为 APS 边界条件依赖于 $\varphi$ 的方式所致。APS 边界条件依赖于 $Y$ 上狄拉克算子的哪些模具有负本征值。对非常小（正）的 $\varrho$，$Y$ 上的狄拉克算子有低能模，而它们的哪些线性组合是负本征值的本征模，取决于 $\varphi$。

% [Tutor Note: 零模附近的“相位转移”机制：$\DD$ 的 $2\times 2$ 块给出 $\lambda=\varrho e^{\i\varphi}$，边界路径积分 $\propto\lambda$。但公式（\ref{durf}）中普法夫量取的是正值 $|\Pf|\propto\varrho$（自然平凡化下），所以相位 $e^{\i\varphi}$ 必须全部由 $\exp(-\i\pi\eta/2)$ 提供——这就是为什么后者在 $\lambda=0$ 处不光滑（其极限依赖于趋近方向 $\varphi$）。APS 边界条件随 $\varphi$ 选择负本征值模，正是这一单值性破缺的来源，也是复费米子反常的“流入”图像。原文第 2262 行 “statment” 系 “statement” 的拼写笔误。]

\subsection{多个能带的情形（The Case of $\nu$ Bands）}\label{bands}

在自由费米子层面，拓扑超导体用一个取整数值的不变量 $\nu$ 来刻画。让我们回忆 $\nu$ 是如何定义的。

在边界上，我们考虑 $n_1$ 个在 $\sT$ 下按
\be\label{trt}\sT\psi(-t,x_1,x_2)=\g_0\psi(t,x_1,x_2),\ee
变换的马约拉纳费米子，以及 $n_2$ 个按相反符号变换的：
\be\label{utrt}\sT\psi(-t,x_1,x_2)=-\g_0\psi(t,x_1,x_2).\ee
一对按相反符号变换的马约拉纳费米子 $\psi,\psi'$ 可以有一个 $\sT$-保持的裸质量 $\bar\psi\psi'$。因此，在 $\sT$-不变微扰的意义下，只有差值 $\nu=n_1-n_2$ 是相关的。

% [Tutor Note: 为何 $\nu=n_1-n_2$ 才是不变量：若 $\psi$ 与 $\psi'$ 在 $\sT$ 下相差一个符号，则双线性型 $\bar\psi\psi'$ 在 $\sT$ 下两个符号因子相互抵消，质量项保持 $\sT$；因此这样的费米子对可以被有隙化而不破坏对称性，只有“符号净数目”有意义。这与第 \ref{bpf} 节体侧的定义一致。]

在体中，如第 \ref{bpf} 节所解释的，基本对象是质量为 $m$、带有质量为 $\mu$ 的正规化子的马约拉纳费米子 $\psi$。有趣的情形是 $m$ 与 $\mu$ 反号。在第 \ref{bpf} 节中，我们取 $\mu>0$、$m<0$，得到 $\exp(-\i\pi\eta/2)$ 作为 sTQFT 配分函数。若符号相反，即 $\mu<0$、$m>0$，同样的推导会给出 $\exp(\i\pi\eta/2)$。在体中，$\nu$ 定义为 $\mu>0,\,m<0$ 的配对数目与 $\mu<0,\,m>0$ 的配对数目之差。体配分函数是 $\exp(-\nu\i\pi\eta/2)$。

在自由费米子层面，$\nu$ 是取整数值的不变量；但已知 \cite{KitTwo}，当计入 $\sT$-保持的相互作用后，$\nu$ 只是 mod 16 意义下的拓扑不变量。

\cite{KTTW} 中指出，对 4 维 $\Pin^+$ 流形，$\exp(-\i\pi\eta/2)$ 总是 \cite{G,Stolz} 1 的 16 次根\footnote{而且所有 16 次根都会出现。对 $X=\RP^4$，有 $\exp(-\i\pi \eta/2)=\exp(\pm 2\i\pi /16)$，其中符号取决于 $\pin^+$ 结构的选择。关于这些事项的入门介绍，见附录 \ref{feta}。}。因此，$\nu$ 个配分函数为 $\exp(-\i\pi\eta/2)$ 的 sTQFT 副本的乘积平凡，当且仅当 $\nu$ 是 16 的整数倍。\cite{KTTW} 以此作为“为何即便在相互作用层面 $\nu$ 也在 mod 16 意义下良定义”的一个解释。我们在第 \ref{bpf} 节中证明与 $D=4$ 拓扑超导体相伴的基本 sTQFT 的配分函数确实就是 $\exp(-\i\pi \eta/2)$，从而把这些考虑建立在了更牢固的基础之上。

% [Tutor Note: 附录 \ref{feta} 给出 $\eta$ 的量子化：$\eta_{\P}+\eta_{\P'}=0 \pmod 4$（式 \ref{simpler}）且 $8\eta_{\P'}=8\eta_{\P}\pmod 4$（式 \ref{impler}），由此 $\eta\in(1/4)\Z$，故 $\exp(-\i\pi\eta/2)$ 是 16 次单位根；对 $X=\RP^4$ 直接算得 $\eta_{\RP^4}=-1/4$，给出 $\exp(\pm 2\i\pi/16)$（符号由 pin+ 结构决定）。“所有 16 次根都会出现”这一点保证 mod 16 分类是饱和的、不能再粗化。这也是本节标题中 $\nu$ 个能带情形与附录 C 论证的呼应之处。]

\subsection{能带数为 8 时的特殊行为（Special Behavior For $\nu=8$）}\label{zeb}

\cite{HCR} 中证明，至少在一个特殊情形下，与映射环面相关的反常探测的是 $\nu$ 模 8 的值，而不是模 16 的值。现在我们可以解释其原因为何如此。首先，$\exp(-\i\pi \eta/2)$ 是 $\pin^+$ 4-流形 $X$ 的一个配边不变量。此外，它生成配边不变量群，所以任何其它 $\U(1)$-值配边不变量都是它的幂。特别地，取 $w_4$ 表示第四 Stiefel--Whitney 类，$(-1)^{w_4}$ 是一个非平凡的 $\U(1)$-值配边不变量，因此它必是 $\exp(-\i\pi \eta/2)$ 的某个幂。这个幂必须是 8，因为 $(-1)^{w_4}$ 取值于 $\{\pm 1\}$：
\be\label{reffo}(-1)^{w_4}=\exp(-8\i\pi \eta/2). \ee
例如，对 $X=\RP^4$，有 $w_4=1$、$\exp(-\pi \i\eta/2)=\exp(\pm 2\i\pi /16)$，符号取决于 $\pin^+$ 结构。

% [Tutor Note: 原文第 2325、2333、2339 行三处 “Stieffel-Whitney” 系 “Stiefel--Whitney” 的拼写笔误，译文按正确拼写处理。式（\ref{reffo}）中幂次 8 不能取别的值：$\exp(-\i\pi\eta/2)$ 是 16 次单位根，而 $(-1)^{w_4}$ 是 2 次单位根，故幂次必为 8 的奇数倍，模 16 唯一。]

于是拓扑超导体的 $\nu=8$ 情形其体配分函数为 $(-1)^{w_4}$。Stiefel--Whitney 类的定义不需要选择自旋或 pin 结构，所以与 $\nu$ 的其它取值的情形不同，$\nu=8$ 的拓扑超导体不依赖于 4-流形 $X$ 的 $\pin^+$ 结构。存在一个完全成熟的、配分函数为 $(-1)^{w_4}$ 的 TQFT。

对于秩 4 的实向量丛——例如 4-流形的切丛——$w_4$ 等于欧拉类的模 2 约化。因此另一种说法是：$\nu=8$ 理论的体配分函数为 $(-1)^{\chi(X)}$，\footnote{我们从 $w_4$ 而非 $\chi$ 出发，是因为像 $w_4$ 这样的 Stiefel--Whitney 数总是配边不变量，这使解释它与 $\exp(-\i\pi \eta/2)$ 的关系更为容易。同样，$\chi$ 与 $w_4$ 的关系是解释 $(-1)^{\chi(X)}$ 为配边不变量最简单的方式。}其中 $\chi(X)$ 是 $X$ 的欧拉示性数。

现在我们可以解释 \cite{HCR} 中发现的结果，并看出它是普遍的。若 $X$ 是映射环面，即它是纤维丛 $X\to S^1$ 的全空间、纤维为 3-流形 $Y$，则 $\chi(X)=0$。证明这一点的一种方式是注意到：一般地，欧拉示性数对纤维化具有乘性。于是 $\chi(X)=\chi(Y)\chi(S^1)$。但该式为零，因为 $\chi(S^1)=0$。

% [Tutor Note: 乘性 $\chi(X)=\chi(Y)\chi(S^1)$ 要求纤维化具有有限型（此处 $Y$ 紧致），对映射环面成立；$\chi(S^1)=0$ 是全部关键。由此 $\nu=8$ 理论在映射环面上配分函数为 $(-1)^0=1$，故映射环面构造完全探测不到 $\nu$ 的 mod 16 信息，只能探测 mod 8——这就是 \cite{HCR} 结果的普遍形式。]

简而言之，我们处于与第 \ref{signp} 节中在可定向流形上讨论同一问题时所发现的类似情形。若 $\nu$ 是 8 的倍数，这保证了 3-流形 $Y$ 上 $\nu$ 个相同马约拉纳费米子的理论没有传统的整体反常——即可通过映射环面的不变量 $\exp(-\nu\i\pi \eta/2)$ 探测到的反常。但若要为 $\nu$ 个相同 $2+1$ 维马约拉纳费米子的系统自然地定义配分函数的符号、而不需要来自 4-流形 $X$ 的“反常流入”，人们需要 $\exp(-\nu\i\pi \eta/2)$ 对所有 4-流形（而不只是映射环面）都等于 1。为此，我们需要 $\nu$ 是 16 的倍数。

% [Tutor Note: 这是“8 与 16”的核心区分：8 只保证映射环面上的相位为 1（传统整体反常消失）；16 才保证对所有 4-流形相位为 1，从而可以不需要体（无需反常流入）而直接定义边界理论的配分函数。两者的差别正是“能否用配边不变量给出一致定义”，而非微扰计算可见的差别。]

$\nu=8$ 的体态具有配分函数 $(-1)^{\chi(X)}$ 这一事实，初看之下似乎提供了一种定义无拓扑序的、保持对称性的有隙边界态的办法。毕竟，带边界流形的欧拉示性数完全良定义，并具有自然的配边性质。在配分函数为 $(-1)^{\chi(X)}$ 的 TQFT 中，人们可以通过规定“任意带空间边界 $Y$ 的 4-流形 $X$ 上的配分函数为 $(-1)^{\chi(X)}$”来定义一个保持对称性的边界态。这个边界态满足 TQFT 的切割与黏合公理，但它与反射正定性或幺正性不相容。举一个简单的反例：取 $Y$ 为三维球面、$X$ 为以 $Y$ 为边界的四维球体。则 $(-1)^{\chi(X)}=-1,$ 但由于 $X$ 可以像第 \ref{rpo} 节的图 \ref{dodotwo} 那样由两个相同的部分黏合而成，反射正定性会要求 $X$ 的配分函数为正。

% [Tutor Note: 反射正定性论证：若 $X$ 由两个相同半体沿边界黏合，则配分函数应为 $|\psi|^2>0$；而四维球体 $\chi=1$ 给出 $(-1)^{\chi(X)}=-1$，矛盾。故 $\nu=8$ 的有隙边界态虽满足切割--黏合，却破坏幺正性——这与第 \ref{overv} 节末段所说的“该边界态与幺正性不相容”是同一事实。]

\subsection{不可定向时空上的三维拓扑绝缘体（The $d=3$ Topological Insulator On An Unorientable Spacetime）}\label{topsupun}

现在我们非常简要地考虑可能不可定向的时空上的 $3+1$ 维拓扑绝缘体。我们关于拓扑超导体所说的一切都有相当明显的类比，只有一个例外。与拓扑超导体的整数值不变量 $\nu$ 相对照，拓扑绝缘体即使在自由费米子层面也只有取值于 $\Z_2$ 的不变量。确实，我们已在式（\ref{worx}）中写下了一对在 $\sT$ 下以相同方式变换的 $2+1$ 维狄拉克费米子 $\psi_1,\psi_2$ 的一个 $\sT$-保持质量项。所以 $\sT$ 对称性只 mod 2 地保护无质量带电狄拉克费米子的数目。

因此我们预期，对拓扑绝缘体而言，即便在不可定向的 4-流形上，$\exp(-\i\pi \eta/2)$ 也总是 $\pm 1$。要理解这一陈述，我们首先必须回忆：保护拓扑绝缘体 SPT 态的对称性，用相对论的术语说，\footnote{$\sC$ 与 $\sCT$ 把电子与正电子互换，它们在凝聚态物理中从来不是真正的对称性。}是 $\sT$ 而不是 $\sCT$。正如我们在第 \ref{welfin} 节所讨论的，这意味着在欧氏号差下我们应当考虑对称性 $\sCR$；换句话说，反转时空定向的对称性同时也起电荷共轭的作用。

% [Tutor Note: 这里的 $\sCR$ 正是“不可定向时空上 $\theta=\pi$”的精确表述所需的结构：在可定向时空上，拓扑绝缘体的 $\theta=\pi$ 项可写作 $\exp(\i\pi P)$（式 \ref{zpart}），其物理解释见关于 $\theta=\pi$ 物理含义的一节；一旦时空不可定向，$P$ 与 $\theta$ 项本身不再良定义，唯一存活的是由 $\sCR$ 保证的 $\Z_2$ 相位 $\exp(-\i\pi\eta/2)=\pm1$。因此不可定向化“丢掉”了 $\theta=\pi$ 的连续信息，只留下一个符号。]

现在设 $X$ 是可能不可定向的 4-流形，$\h X$ 是其可定向双重覆盖。我们想在 $X$ 上理解的 $\exp(-\i\pi \eta/2)$ 是作用在费米子场 $\psi^\pm _\pm$ 上的狄拉克算子的 $\eta$ 不变量，其中上标是 $\U(1)$ 荷、下标是手征性。在 $\h X$ 上有一个等价定义：只考虑具有正 $\U(1)$ 荷但任意手征性的费米子场 $\h\psi^+_\pm$；由仅作用于 $\h\psi^+_\pm$ 的狄拉克算子的本征值算出的 $\h X$ 上狄拉克谱，与由 $X$ 上全体场 $\psi^\pm_\pm$ 算出的狄拉克谱相同。确实，$X$ 是 $\h X$ 按 $\sCR$ 对称性取的商，所以 $X$ 上的场 $\psi^+_+$ 与 $\psi^-_-$ 合为 $\h X$ 上的 $\h\psi^+_+$，类似地 $X$ 上的 $\psi^+_-$ 与 $\psi^-_+$ 合为 $\h X$ 上的 $\h\psi^+_-$。因此，不必通过对 $X$ 上作用在全体场 $\psi^\pm_\pm$ 上的狄拉克算子的谱求和来计算 $\eta$ 不变量，我们可以改为对 $\h X$ 上仅作用在 $\psi^+_\pm$ 上的狄拉克算子的谱求和。

% [Tutor Note: 这是本节最技术性的一步：$\sCR$ 同时反转定向并做电荷共轭，故在 $\h X$ 上可以把 $X$ 的手征场两两配对成单一正荷场 $\h\psi^+_\pm$，手征性在 $\h X$ 上良定义；于是不可定向 $X$ 上的 $\eta$（对 $\psi^\pm_\pm$ 全体）等于可定向 $\h X$ 上的 $\eta$（只对 $\h\psi^+_\pm$）。这使下面的谱对称性论证成为可能。]

由于 $\h X$ 可定向，$\h X$ 上的狄拉克谱在 $\lambda\leftrightarrow -\lambda$ 下对称。这照例由手征性的考虑推出。但一对本征值 $\lambda,-\lambda$ 对 $\eta$ 不变量没有贡献。所以在 $\h X$ 上计算 $\eta$ 时，我们可以只考虑零模。$\h X$ 上狄拉克算子的零模空间由于狄拉克方程通常的反线性对称而是偶维的，因此 $\eta$ 是偶整数。因此，$\exp(-\i\pi \eta/2)=\pm 1$，这正是我们想要证明的。

% [Tutor Note: 零模空间偶维来自狄拉克方程的反线性对称（在 4 维 $\Pin^+$ 情形即 Kramers 型对称），故 $\eta\in 2\Z$，$\exp(-\i\pi\eta/2)=(-1)^{\eta/2}=\pm1$；这正是拓扑绝缘体 $\Z_2$ 不变量的场论表述，与拓扑超导体的 $\nu\in\Z$（16 次根）形成对照。此处只用到谱对称性与零模计数，不需要显式计算 $\eta$，这是原文省略的推理环节。]

\section{马约拉纳链（The Majorana Chain）}\label{maj}

\begin{keytakeaways}
\begin{itemize}[leftmargin=1.6em,itemsep=0.2em,topsep=0.2em]
  \item \textbf{物理图像}：1 维马约拉纳链是时间反演平方为 1 的对称性所保护的 SPT 相；其 0 维边界上的实费米子正是端点马约拉纳零模的低能有效描述，边界零模的有无给出 Z2 指标，相互作用下只保留模 8 的分类。
  \item \textbf{数学工具}：0+1 维与 1+1 维的 pin 结构、欧氏号差下的赝实表示、二维 APS 不变量的半整数量子化（8 次单位根），以及 Dai--Freed 定理。
  \item \textbf{核心结论}：时间反演平方为 1 允许保持对称性的裸质量，故二维马约拉纳费米子在欧氏号差下是赝实的；体自旋拓扑量子场论的配分函数相位是 8 次单位根，能带数为 8 的倍数时体理论平凡；边界反常与体相位由同一个 APS 不变量统一刻画，并与附录 C 关于四维 APS 不变量的论证相呼应。
\end{itemize}
\end{keytakeaways}

在微观层面，现实世界中成立的时间反演操作满足 $\sT^2=(-1)^F$，其中 $(-1)^F$ 是模 2 数费米子的算子。但在凝聚态物理中有一些情形，考虑满足 $\sT^2=1$ 的时间反演对称是自然的。这通常发生在两种情况下：要么自旋不重要，要么所考虑的对称性实际上是时间反演与一个空间坐标反射的乘积 $\sRT$。（在自然界中 $(\sRT)^2=1$，而 $\sRT$ 可以作为时间反演作用在局域于 $\sR$ 的不动点集上的薄层或长链上。）

在物质 SPT 相的语境中，有一个非常有趣的例子，其中假定了满足 $\sT^2=1$ 的时间反演对称性。这就是 1 维空间中的马约拉纳链 \cite{FK}。

这类系统的边界是 0 维的，也就是一个点，我们将从那里开始。在本文中，我们通过取洛伦兹号差下的伽马矩阵为满足 $\{\g_\mu,\g_\nu\}=2\eta_{\mu\nu}$ 的实矩阵（其中 $\eta_{\mu\nu}=\diag(-1,1,1\dots,1)$）把条件 $\sT^2=(-1)^F$ 内置进来。然后我们取 $\sT$ 的作用为
\be\label{horz}\psi(t,\vec x)
\to \pm \g_0\psi(-t,\vec x),\ee
从而保证 $\sT^2\psi=-\g_0^2\psi=-\psi$。由于对 $i>0$ 有 $\g_i^2=+1$，由此可知空间反射的平方为 $+1$。在欧氏号差下，这导致 $\pin^+$ 结构。

% [Tutor Note: 原文此处疑有符号笔误（第 2420 行）：由 $\{\g_\mu,\g_\nu\}=2\eta_{\mu\nu}$、$\eta_{00}=-1$ 得 $\g_0^2=-1$；而把 $\sT\psi(t,\vec x)=\pm\g_0\psi(-t,\vec x)$ 作用两次（$\sT$ 反线性、$\g_0$ 为实矩阵）直接得到 $\sT^2\psi=\g_0^2\psi=-\psi$，并没有额外的负号。原文写作 $-\g_0^2\psi$ 与 $\g_0^2=-1$ 及结论 $\sT^2=(-1)^F$ 不相容。译文保留原式以便对照，读者按 $\sT^2\psi=\g_0^2\psi=-\psi$ 理解即可。]

在 $0+1$ 维中，我们只需要一个伽马矩阵 $\g_0$。为得到尽可能经济的理论，我们取 $\g_0^2=+1$，而不是像此前所有讨论中那样取 $-1$。于是我们可以用 $1\times 1$ 矩阵 1（或 $-1$）表示 $\g_0$，并且可以考虑一个单分量实（马约拉纳）费米子场 $\psi$，其作用量为
\be\label{dolb}I=\frac{i}{2}\int \d t \,\psi\frac{\d}{\d t}\psi. \ee
我们取 $\sT$ 作用在这样的场上为
\be\label{olb}\sT\psi(t)=\pm \psi(-t), \ee
并作某种符号选择，使得 $\sT^2=1$。一对按相反符号变换的费米子 $\psi$、$\psi'$ 可以有一个 $\sT$-不变的“质量”项 $i\psi\psi'$，因此如第 \ref{bands} 节那样，自然的不变量是按 $\sT$ 以 $+$ 号变换的费米子数目与以 $-$ 号变换的费米子数目之差 $\nu$。在自由费米子层面，$\nu$ 是整数值不变量；但在相互作用层面，\cite{FK} 指出 $\nu$ 只是 mod 8 意义下的不变量。（这句话的意思是：由 8 个全部以同一符号在 $\sT$ 下变换的实费米子组成的 $0+1$ 维系统，可以被适当的 $\sT$-不变四费米子耦合有隙化。）

% [Tutor Note: 这是马约拉纳链的核心图像：$0+1$ 维单个实费米子 $\psi$ 的希尔伯特空间是 2 维（$|0\rangle$ 与 $\psi|0\rangle$），它正是 1 维马约拉纳链端点零模的低能有效描述；有隙相与平庸相的区分对应这个 2 维空间是否被 $\sT$-不变微扰解除，即 $\Z_2$ 指标。8 个这样的边界费米子可被四费米子相互作用有隙化（Fidkowski--Kitaev），故相互作用下只剩 mod 8 分类——与第 \ref{bands} 节的 mod 16 形成维数上的对照。]

当 $\nu$ 不是 8 的倍数时，这个系统在对称性 $\sT$ 与 $(-1)^F$ 的实现上有一个微妙的（依赖于 $\nu$ 的）反常。我们在此不再重复这个故事，转而要问：带有 $\nu$ 个边界费米子的系统，如何能自然地出现在 $1+1$ 维的 $\sT$-保持系统的边界上。我们在建立这个问题时所作的选择 $\g_0^2=+1$ 对应于 $\Pin^-$ 结构；而一旦我们在时空的 $0+1$ 维边界上假定了 $\Pin^-$ 结构，我们就必须（如果希望 Dai--Freed 定理或类似工具可用）在 $1+1$ 维体中也假定 $\Pin^-$ 结构。因此我们考虑 $1+1$ 维的伽马矩阵满足 $\g_0^2=1=-\g_1^2$、或更协变地满足 $\{\g_\mu,\g_\nu\}=-2\eta_{\mu\nu}$ 的费米子。在欧氏号差下，这变成 $\{\g_\mu,\g_\nu\}=-2\delta_{\mu\nu}$，其中负号对应于 $\Pin^-$ 结构。（满足 $\{\g_\mu,\g_\nu\}=-2\delta_{\mu\nu}$ 的 $2\times 2$ 矩阵存在，但不能取为实矩阵。如下一段所讨论的，裸质量是可能的这一事实清楚表明：具有 $\sT^2=1$ 的马约拉纳费米子在欧氏号差下是赝实的。）

% [Tutor Note: pin 结构的对应关系是全节的枢纽：边界 $0+1$ 维取 $\g_0^2=+1$（使 $\sT^2=1$）对应 $\Pin^-$，为让 Dai--Freed 定理适用，$1+1$ 维体也必须是 $\Pin^-$（$\{\g_\mu,\g_\nu\}=-2\delta_{\mu\nu}$）。欧氏 2 维中该代数存在但不能全取实矩阵，正说明 $\sT^2=1$ 的马约拉纳费米子在欧氏号差下是赝实而非实——这与它能带 $\sT$-保持裸质量互为表里。]

与 $\sT^2=-1$ 的情形相对照，具有 $\sT^2=1$ 的 2 维马约拉纳费米子可以有一个 $\sT$-不变的裸质量。在洛伦兹号差下，$\sT$-不变的狄拉克方程可写成
\begin{equation}\label{sx}\left(\g^0\partial_0+\g^1\partial_1-m\bar\g\right)\psi=0, \ee
其中 $m$ 为实数，$\bar \g=\g^0\g^1$ 是手征算子。其实，该方程对 $\{\g_\mu,\g_\nu\}=\pm 2\eta_{\mu\nu}$ 中的任一符号都是 $\sT$-不变的；但若把方程写成某个线性算子 $H$ 的哈密顿形式 $\i\partial_t\psi=H\psi$，就会看到只有当符号取 $-$ 时 $H$ 才是厄米的。所以正是那种情形下这个质量项在物理上才是合理的。

% [Tutor Note: 原文此处把环境写成 \begin{equation} 配 \ee（第 2452 行），混用了环境与宏；因导言区将 \ee 定义为 \end{equation}，编译无误，译文保留原样。物理要点：$\bar\g=\g^0\g^1$ 在此 2 维中满足 $\bar\g^2=1$；只有 $\{\g_\mu,\g_\nu\}=-2\eta_{\mu\nu}$ 时哈密顿量 $H$ 厄米，故 $\Pin^-$ 才是质量项有物理意义的情形。]

推广到弯曲时空并延拓到欧氏号差后，式（\ref{sx}）变成 $\left(\slashed{D}-\i m\bg \right)\psi=0$。但乘以 $\bg$ 并把方程写成
\be\label{mirth}\left(\D'+\i m\right)\psi=0.  \ee
更为方便。这里 $\D'=\bg \slashed{D}$ 是一个厄米且反射对称的狄拉克算子。由于它反射对称，它在任何赋有 $\Pin^-$ 结构的 2-流形上（可能不可定向）都有意义。（关于这一点，见附录 \ref{do}。）

马约拉纳链的体 sTQFT 由如下方式构造：取一个质量为 $m$ 的费米子场 $\psi$，连同质量为 $\mu$ 的泡利--维拉尔斯正规化子场。如第 \ref{bpf} 节那样，有趣的情形是 $m$ 与 $\mu$ 反号。取 $\mu>0$、$m<0$，由熟悉的论证，配分函数的相位为 $\exp(-\i\pi \eta/2)$。这就是与马约拉纳链相伴的体 sTQFT 的配分函数。它总是 1 的 8 次根；对 $\RP^2$ 它等于 $\exp(\pm 2\i\pi /8)$，符号取决于 $\Pin^-$ 结构的选择。（见附录 \ref{feta}。）因为 $\exp(-\pi \i\eta/2)$ 总是 1 的 8 次根，马约拉纳链的体 sTQFT 对 $\nu$ 为 8 的倍数时是平凡的。（\cite{KTTW} 基于配边不变性的假定论证过这一点。）

% [Tutor Note: 8 次根与 16 次根的对照：$D=2$ 的 $\Pin^-$ 情形中 $\eta\in(1/2)\Z$（附录 \ref{feta}：$\eta_{\RP^2}=\pm1/2$），故 $\exp(-\i\pi\eta/2)$ 是 8 次单位根，对应相互作用下的 mod 8 分类；而第 \ref{bands} 节的 $D=4$、$\Pin^+$ 情形 $\eta\in(1/4)\Z$ 给出 16 次根与 mod 16。维数降低一维，$\eta$ 的分母减半，分类模数也随之减半——这是马约拉纳链与 3 维拓扑超导体分类差异的几何根源。]

一如既往，若 $\nu$ 不是 8 的倍数，体配分函数 $\exp(-\nu\i\pi \eta/2)$ 无法在带边界的、可能不可定向的 2-流形上以保持对称性的方式定义为拓扑不变量。然而，Dai--Freed 定理适用于这种情形，它指出：$\nu$ 个边界费米子的路径积分与 $\exp(-\nu\i\pi \eta/2)$ 的乘积是良定义的。这就是由体 sTQFT 与边界费米子组成的复合系统的路径积分。

% [Tutor Note: 这是第 \ref{comb} 节 Dai--Freed 结论在 $D=2$、$\Pin^-$ 情形下的重述，请与式（\ref{durf}）及附录 \ref{feta} 对读：$\eta$ 的两种取值对应 $\RP^2$ 上两种 $\Pin^-$ 结构的选择，正如 4 维 $\RP^4$ 上两种 $\pin^+$ 结构。至此，§4 的 16 次根与 §5 的 8 次根、以及“$\nu$ 个副本何时平凡”的判据，都由同一个 $\eta$ 不变量框架统一说明。]
